% ============================================================================
%  \section{System Model}  —  PAPER VERSION, IEEEtran conference
%  Drafted 2026-09-07.  Designed for the 15 Sep paper.
%
%  MEASURED: ~1500 words prose + 4 display equations = ~2.7 columns.  With
%  Related Work at 2.4 columns that is too much for a 6-page paper.  To reach
%  ~1.8 columns make these cuts, in order (each self-contained):
%     1. §Parameter Studies -> two sentences; it duplicates §Experiments (-170 w)
%     2. §Generative, 2nd para (the raw-IQ / mode-collapse motivation) ->
%        move to Methodology, where it belongs                        (-130 w)
%     3. §Attacker Ladder, 2nd para (the "boundary is the bar" argument) ->
%        move to Methodology or Related Work                          (-110 w)
%  Do NOT cut: the two-views paragraph in §Link (it sets up the whole
%  effectiveness/detectability split), the D_NP justification in §Defender
%  (it is why the model was reduced at all), or the first paragraph of
%  §Cooperative (it is the correction to the coordination hypothesis).
%
%  Replaces main.tex §System Model, which still describes the K-subcarrier /
%  TDL / N_J-jammer setting that no working experiment supports (README
%  "Thesis proposal + settled framing", 2026-09-01).
%
%  Satisfies, by construction:
%    SUPERVISOR_TODO §3  where detection happens, what the detector observes,
%                        assumption tiers, coordination assumption, CTDE split,
%                        counter-signal non-viability, localization out of scope
%    SUPERVISOR_TODO §4  reward = BER - beta*det, power as HARD constraint
%    SUPERVISOR_TODO §8  the five-baseline envelope
%    SUPERVISOR_TODO §11 M0: single subcarrier, one channel, NP-optimal test
%    SUPERVISOR_TODO §13 sigma as primary swept axis, sigma=0 anchor
%
%  Style per A. Di Maio: one sentence per line; \cite before all punctuation
%  with a ~ tie; text reads correctly with the brackets removed.
%
%  NEW KEYS NEEDED IN refs.bib (see end of file for stubs):
%      zhou2025cgan, valianti2024cooperative (exists), amuru2015optimal (exists)
% ============================================================================

\section{System Model}
\label{sec:system}

\subsection{Link and Signal Model}
\label{sec:system:link}

We model a single-carrier link carrying $N$ QPSK symbols per frame.
The victim transmits $s[n] \in \mathcal{S} = \{(\pm 1 \pm j)/\sqrt{2}\}$, drawn independently and uniformly, reaching the receiver through a channel coefficient $h_0 \in \mathbb{C}$.
A team of $K$ jammers indexed by $k \in \mathcal{K}$ transmits $x_k[n]$, reaching the same receiver through $g_k \in \mathbb{C}$.
The receiver observes
\begin{equation}
  y[n] \;=\; h_0\, s[n] \;+\; \sum_{k \in \mathcal{K}} g_k\, x_k[n] \;+\; w[n],
  \label{eq:rx}
\end{equation}
with $w[n] \sim \mathcal{CN}(0, 2\sigma^2)$ and $\sigma$ the noise standard deviation per real dimension, so that $N_0 = 2\sigma^2$.
Equalising by $h_0$ gives the decision variable
\begin{equation}
  \hat{y}[n] \;=\; s[n] \;+\; d[n] \;+\; \sigma_{\mathrm{eff}} \tilde{w}[n],
  \qquad
  d[n] \;\triangleq\; \sum_{k \in \mathcal{K}} \underbrace{\frac{g_k}{h_0} x_k[n]}_{\;\triangleq\; d_k[n]},
  \label{eq:eq}
\end{equation}
where $d[n]$ is the aggregate perturbation in the equalised domain and $d_k[n]$ is agent $k$'s effective contribution to it.
The single-jammer, unit-channel case $K=1$, $h_0 = g_1 = 1$ is the base model; every extension in Sec.~\ref{sec:system:ladder} reintroduces one term of \eqref{eq:rx}.

Two views of the same frame carry the whole tension of the problem.
The victim decides on $\hat{y}$ and is therefore harmed by $d$ alone, whereas the detector observes the composite $y$ before equalisation, so it sees the interference through $g_k$ rather than through $g_k / h_0$.
Effectiveness is a property of the equalised perturbation and detectability is a property of the received composite, and the two are related but not identical --- which is the design margin the attacker exploits.

Each jammer is constrained by a hard per-frame power budget
\begin{equation}
  \frac{1}{N} \sum_{n=1}^{N} \bigl| x_k[n] \bigr|^2 \;\le\; P_k ,
  \label{eq:power}
\end{equation}
imposed as a projection onto the feasible set at the output of the attacker rather than as a penalty term, so that no attacker in this paper can buy effectiveness by paying a soft power cost.

Performance at the victim is the bit error rate and the symbol error rate under maximum-likelihood QPSK detection.
Bit-error recovery is out of scope: we assume forward error correction, retransmission and rate adaptation are handled by a higher layer, and report the raw error rate because it is the input any such layer receives, and because pushing it past the code's correcting capability is what becomes outage.

\subsection{Attacker Model}
\label{sec:system:attacker}

We model every attacker, learned or not, as a conditional law
\begin{equation}
  d_k[n] \;\sim\; \pi_k\bigl( \cdot \mid c_k[n] \bigr),
  \label{eq:policy}
\end{equation}
where $c_k$ is the information available to agent $k$ when it forms the perturbation.
This is the paper's central abstraction, and it is what places the closed-form and the learned attacks on a single axis: the strategies differ only in the contents of $c_k$ and in the family of laws $\pi_k$ ranges over.

Three assumption tiers define $c_k$, and results are reported per tier so that no capability is granted silently.
At tier~T0 the attacker knows the victim symbol $s[n]$, both channel coefficients $h_0$ and $g_k$, and is sample-synchronous with the preamble; this is a genie and serves only as an upper bound.
At tier~T1 the attacker holds a channel estimate $\hat{g}_k$ and achieves preamble synchronisation with residual carrier-frequency offset, timing offset and phase error, and does not know $s[n]$.
At tier~T2 it holds neither.
The counter-signal attack exists only at T0, and its non-viability at T1 is structural rather than incidental: cancellation is an exact-inverse operation, so residual phase error destroys it, whereas an attack that only has to place the received point in the wrong half-plane degrades gracefully.

Coordination is assumed to be a shared clock and a command codebook agreed offline, with no per-frame backhaul between jammers.
This assumption is load-bearing and is chosen deliberately: a per-frame backhaul would let the team behave as one jammer with $K$ antennas, which both trivialises the coordination question and is the least realistic hardware assumption available.
Under centralised training with decentralised execution~\cite{lowe2017maddpg}, the error rate is available to a centralised critic during training and is \emph{not} available to a deployed jammer, whose execution-time observation is its own transmission, its own channel estimate, and at most a delayed acknowledgement from the victim.
The objective below is therefore a training-time construct, and we label it as such; it is the attacker-side mirror of the defender's inability to obtain ground-truth labels at run time.

\subsection{Defender Model}
\label{sec:system:defender}

Detection takes place at the victim receiver, on the composite pre-equalisation frame $y$ of \eqref{eq:rx}.
The detector observes that frame and nothing else: it has no channel state information, no side information about the jammers, and no ground-truth labels at execution time, so it cannot adapt online and any adaptation is round-based and offline.
We evaluate against a suite of three detectors at a common false-alarm rate $\alpha$, calibrated on clean frames at the operating noise level.

The energy detector $D_{\mathrm{E}}$ thresholds the mean frame power $\frac{1}{N}\sum_n |y[n]|^2$, and is the classical baseline any added-power attack must survive~\cite{urkowitz1967energy}.
The learned detector $D_{\mathrm{L}}$ is a small convolutional classifier on the two-dimensional in-phase/quadrature histogram of $y$, the single-carrier counterpart of the spectrogram classifiers that define the state of the art in this setting~\cite{li2022jamming_ofdm}.
The reference detector $D_{\mathrm{NP}}$ is the Neyman--Pearson likelihood-ratio test, which the single-carrier model makes available in closed form: with $\sigma$ known and the perturbation law $\pi$ stated,
\begin{equation}
  \Lambda(y) = \sum_{n=1}^{N} \log \frac{p_1\bigl(y[n]\bigr)}{p_0\bigl(y[n]\bigr)},
  \quad
  p_0(y) = \tfrac{1}{4}\!\sum_{s \in \mathcal{S}}\! \mathcal{CN}\bigl(y; h_0 s, 2\sigma^2\bigr),
  \label{eq:np}
\end{equation}
and $p_1$ is the same mixture with $h_0 s$ replaced by $h_0(s + d)$ and averaged over $d \sim \pi$.

$D_{\mathrm{NP}}$ is the reason the model is worth reducing to a single carrier.
It converts every claim of the form ``the trained detector does not see this attack'' into the far stronger ``no detector can do better than this'', and it supplies the reference point that an adaptation-cost measurement otherwise lacks, since the distance between a retrained detector and the optimal one is precisely the adaptation budget that remains.

Jammer localisation --- angle-of-arrival, time-difference-of-arrival, or received-signal-strength anomaly --- is a real defender capability and is explicitly out of scope, since it requires multi-antenna or distributed infrastructure that the baseband model does not represent.
We note it as future work rather than assume it away.

\subsection{Objective and Reporting Convention}
\label{sec:system:objective}

The attacker maximises the induced error rate penalised by detection,
\begin{equation}
  \max_{\{\pi_k\}} \;\; \mathrm{BER}\bigl(\pi_{1:K}\bigr) \;-\; \beta \, P_{\mathrm{det}}\bigl(\pi_{1:K}; \alpha\bigr)
  \quad \text{s.t.} \quad \eqref{eq:power},
  \label{eq:obj}
\end{equation}
with no further terms: the power budget is a constraint of the environment and every other consideration is a property of the environment rather than of the reward.

Two reporting conventions are non-negotiable and both correct errors this line of work has already made.
First, strategies are compared at \emph{matched detectability} rather than at matched transmit power or matched configuration, since an attack that raises the error rate by raising its own signature is not an improvement.
Second, the stealth budget is set to the defender's own clean false-alarm rate rather than to a conventional threshold such as $P_{\mathrm{det}} \le 0.5$, because an attacker caught in half of every frame is caught within a few frames.
Results are reported as the achievable frontier $\mathrm{BER}^\star(\beta) = \max\{\mathrm{BER} : P_{\mathrm{det}} \le \beta\}$, together with the dual-axis view of error rate and detection probability against the swept parameter, with the no-attacker floor and the omniscient ceiling drawn in.

\subsection{Attacker Ladder}
\label{sec:system:ladder}

Five strategies instantiate \eqref{eq:policy}, and all five appear in every results figure so that any curve can be read against both its floor and its ceiling.
\emph{No attacker} fixes the error-rate floor and the detector's false-alarm rate, the lower reference on both axes.
\emph{Barrage} draws $d[n]$ isotropically at the power budget, knows nothing, and is the naive floor.
The \emph{minimum-energy boundary} attack is closed form and requires no training: for unit-energy QPSK it displaces each symbol just across its nearest decision boundary, $d[n] = -\rho\,\Re\{s[n]\}$ or its quadrature counterpart, pre-compensated as $x_k[n] = d[n] h_0 / g_k$, with $\rho$ sweeping the frontier~\cite{amuru2015optimal}.
The \emph{counter signal} $x_k = -2 h_0 s / g_k$ is the omniscient ceiling requested as an impossible-to-beat reference, and is reported as a bound rather than claimed as an attack.
The \emph{generative} attacker, single-agent and then cooperative, is the proposed method.

The boundary attack is the bar that matters.
Against an independent and uniformly distributed payload the minimum-energy perturbation is available analytically, so on a single link a learner can at best rediscover it; the learned attacker must therefore be motivated by structure that survives scrambling, of which the relevant kinds are the detector's own decision surface, the per-link gains and phases, and the coordination problem itself.

\subsection{Generative Attack Policy}
\label{sec:system:generator}

We parameterise $\pi_k$ by a conditional generator $G_{\theta_k}(z; c_k) \mapsto d_k$ with $z \sim \mathcal{CN}(0,\mathbf{I})$, followed by the projection enforcing \eqref{eq:power}.
Conditional generation has been applied to jamming waveform synthesis before~\cite{zhou2025cgan}, and two departures from that use are what make it an evasion attack rather than an effectiveness attack.
The discriminator is the \emph{detector} operating on the received composite $y$, not a similarity model separating the interference from the victim's signal, so the generator is rewarded for making the received frame indistinguishable from an un-jammed one rather than for making its own emission resemble the victim's.
The objective carries both terms of \eqref{eq:obj}, whereas prior generative jamming optimises effectiveness alone and reports it against a theoretical jamming bound with no detector present.

The generator also resolves a structural obstacle rather than adding capability for its own sake.
A policy acting directly on the transmitted samples has an action space of $2N$ real dimensions driven by a frame-level scalar reward whose variance across a training batch is close to zero, which leaves a policy-gradient method with no usable gradient; parameterising the perturbation \emph{law} instead reduces the coordinating policy's action to a low-dimensional command $m_k$ that conditions a separately trained generator.
Conditioning also pins the generated mode, which is the specific answer to the mode-collapse objection that governs the choice between adversarial and flow-based generators in this setting.

\subsection{Cooperative Attack}
\label{sec:system:coop}

The detector observes the sum $\sum_k g_k x_k$, so any two power allocations producing the same sum are the same observation and splitting power across agents is, by itself, worth nothing at the victim's detector.
Cooperation must therefore act on the \emph{law} of the aggregate perturbation rather than on its magnitude, and two mechanisms do so.

First, at fixed aggregate received perturbation power, independent per-agent generators make $d = \sum_k d_k$ approach a Gaussian law as $K$ grows.
A Gaussian perturbation is distinguishable from an increase in the noise floor only through its energy, so the optimal test of \eqref{eq:np} degenerates towards the energy detector and the team converts structure into stealth at constant effect.
Whether that conversion is free is the question this paper puts: a boundary-directed perturbation is a function of $s[n]$ and is therefore correlated across any agents that share it, so the team must trade the energy efficiency of a directed attack against the stealth of an undirected one, and where that trade settles is not available in closed form.

Second, with $V > 1$ legitimate receivers the team maximises aggregate error subject to $\max_v P_{\mathrm{det}}^{(v)} \le \beta$, and the concavity of the error rate in perturbation power makes spreading strictly preferable to concentration at matched worst-case detectability.
This is the setting in which the per-link geometry $g_{k,v}$ does work, and it is where the cooperative jamming literature's coordination structure transfers~\cite{valianti2024cooperative}: an inter-agent coupling constraint that renders the team problem non-separable, with the mutual-interference limit of that work replaced here by a joint detectability limit.

\subsection{Parameter Studies}
\label{sec:system:params}

The noise level is the primary swept axis, on a logarithmic grid over $\sigma \in [10^{-3}, 0.5]$ with $\sigma = 0$ as a separate anchor.
The anchor is degenerate and is included for that reason: with no noise the un-jammed constellation is four exact points, any perturbation is detected with probability one, and the comparison collapses to error rate at matched power.
Stealth is thus a consequence of noise, and stating so in a model where it can be derived is the point of the reduction.
Secondary axes are the power budget, the number of jammers, the number of legitimate receivers, the attacker synchronisation quality, and the exploitable structure in the transmitted sequence.

% ----------------------------------------------------------------------------
% refs.bib stub for the one new key:
%
% @inproceedings{zhou2025cgan,
%   author    = {Zhou, Liangsi and Tan, Yihan and Xu, Hua},
%   title     = {Communication Jamming Waveform Generation Technology Based on
%                Conditional Generative Adversarial Networks},
%   booktitle = {Proc. 4th Int. Symp. Semiconductor and Electronic Technology (ISSET)},
%   address   = {Xi'an, China},
%   pages     = {373--377},
%   year      = {2025},
%   doi       = {10.1109/ISSET66828.2025.11184988}
% }
% ----------------------------------------------------------------------------
